% ---------------------------------------------------------------------
%  CHƯƠNG 3 — PHÂN TÍCH VÀ THIẾT KẾ
% ---------------------------------------------------------------------
\chapter{Phân tích và thiết kế}
\label{ch:thietke}

\begin{muctieu}
Chương này phân tích đầu vào, đầu ra và yêu cầu của hệ thống; đồng thời mô
tả kiến trúc, dữ liệu và cách tổ chức các thành phần nhận diện, gợi ý phối
đồ và thử trang phục ảo.
\end{muctieu}

\section{Phân tích bài toán và yêu cầu}

Hệ thống hướng đến người dùng Android có nhu cầu lưu trữ quần áo cá nhân và
nhận hỗ trợ khi lựa chọn trang phục. Đầu vào chính là ảnh có một trang phục
nổi bật, ảnh người dùng phục vụ thử đồ và thông tin tài khoản. Đầu ra gồm
loại quần áo dự đoán, độ tin cậy, tối đa ba phương án phối đồ và ảnh xem
trước sau khi ghép lớp trang phục.

Các yêu cầu chức năng được xác định từ phạm vi ở Chương~\ref{ch:modau} gồm:
\begin{itemize}[leftmargin=1.6em, itemsep=3pt]
  \item đăng ký, đăng nhập và quản lý thông tin người dùng;
  \item nhận ảnh từ camera hoặc thư viện ảnh;
  \item nhận diện một số nhóm quần áo và trả độ tin cậy;
  \item thêm, sửa, xoá và xem các món đồ trong tủ đồ cá nhân;
  \item gợi ý tối đa ba phối đồ theo loại, màu sắc và phong cách;
  \item lưu trang phục, phối đồ hoặc kết quả yêu thích;
  \item ghép ảnh quần áo nền trong suốt lên ảnh người dùng và điều chỉnh vị
        trí, kích thước, góc xoay.
\end{itemize}

Yêu cầu phi chức năng gồm giao diện sử dụng được trên thiết bị Android, dữ
liệu mỗi người dùng được tách biệt, lỗi API được xử lý rõ ràng và thời gian
phản hồi hướng tới không quá 3 giây trong điều kiện kiểm thử ổn định.

Phạm vi tương tác giữa người dùng và các chức năng chính được tổng hợp ở
Hình~\ref{fig:usecase-tongquat}. Sơ đồ tập trung vào một tác nhân người dùng
và bảy use case đã xác định từ phạm vi đồ án; quan hệ với API và Firebase
được trình bày riêng trong sơ đồ kiến trúc ở Hình~\ref{fig:kientruc}.

\begin{figure}[H]
\centering
\resizebox{0.94\textwidth}{!}{%
\begin{tikzpicture}[
  uc/.style={draw=daumaroon, ellipse, align=center, fill=lightwarm,
             minimum width=3.5cm, minimum height=1.05cm, font=\small},
  assoc/.style={draw=gray700, line width=0.55pt}
]
  \draw[draw=daunavy, rounded corners=4pt, line width=0.8pt]
    (2.2,0.9) rectangle (12.0,-7.8);
  \node[anchor=west, font=\bfseries\color{daunavy}] at (2.45,0.55)
    {Hệ thống gợi ý phối đồ};

  % Tác nhân người dùng
  \draw[daunavy, line width=0.8pt] (0.7,-2.0) circle (0.28);
  \draw[daunavy, line width=0.8pt] (0.7,-2.28) -- (0.7,-3.25);
  \draw[daunavy, line width=0.8pt] (0.15,-2.62) -- (1.25,-2.62);
  \draw[daunavy, line width=0.8pt] (0.7,-3.25) -- (0.2,-4.0);
  \draw[daunavy, line width=0.8pt] (0.7,-3.25) -- (1.2,-4.0);
  \node[align=center, font=\small\bfseries, text=daunavy] at (0.7,-4.45)
    {Người dùng};

  \node[uc] (auth) at (4.5,-0.8) {Đăng ký, đăng nhập};
  \node[uc] (capture) at (9.3,-0.8) {Chọn hoặc chụp ảnh};
  \node[uc] (wardrobe) at (4.5,-2.65) {Quản lý tủ đồ};
  \node[uc] (recognize) at (9.3,-2.65) {Nhận diện trang phục};
  \node[uc] (recommend) at (4.5,-4.5) {Nhận gợi ý phối đồ};
  \node[uc] (favorite) at (9.3,-4.5) {Lưu và yêu thích};
  \node[uc] (tryon) at (6.9,-6.55) {Thử trang phục 2D};

  % Các đoạn trùng nhau được gom theo hai trục để tránh cắt qua ellipse.
  \draw[assoc] (1.25,-2.62) -- (2.65,-2.62);
  \draw[assoc] (2.65,-0.8) -- (2.65,-6.55);
  \draw[assoc] (2.65,-0.8) -- (auth.west);
  \draw[assoc] (2.65,-2.65) -- (wardrobe.west);
  \draw[assoc] (2.65,-4.5) -- (recommend.west);
  \draw[assoc] (2.65,-6.55) -- (tryon.west);

  \draw[assoc] (1.25,-2.62) -- (1.65,-2.62) -- (1.65,-7.45)
    -- (11.55,-7.45) -- (11.55,-0.8);
  \draw[assoc] (capture.east) -- (11.55,-0.8);
  \draw[assoc] (recognize.east) -- (11.55,-2.65);
  \draw[assoc] (favorite.east) -- (11.55,-4.5);
  \draw[assoc] (6.9,-7.45) -- (tryon.south);
\end{tikzpicture}}
\caption{Sơ đồ use case tổng quát của hệ thống gợi ý phối đồ.}
\label{fig:usecase-tongquat}
\end{figure}

Các use case trong sơ đồ được đặc tả ngắn gọn ở
Bảng~\ref{tab:usecase}. Đầu vào và kết quả chỉ mô tả những gì đã được xác
định trong phạm vi đồ án, không giả định tên màn hình hoặc endpoint cụ thể.

\begin{table}[H]
\centering
\caption{Danh sách use case chính của hệ thống.}
\label{tab:usecase}
\footnotesize
\begin{tabularx}{\textwidth}{>{\bfseries\color{daumaroon}}l >{\bfseries\color{daunavy}\raggedright\arraybackslash}p{3.0cm} >{\raggedright\arraybackslash}X >{\raggedright\arraybackslash}X}
\toprule
\rowcolor{gray100}\bfseries\color{daumaroon}Mã & \bfseries\color{daunavy}Use case & \bfseries\color{daumaroon}Đầu vào/thao tác & \bfseries\color{daunavy}Kết quả mong đợi \\
\midrule
UC01 & Xác thực & đăng ký hoặc đăng nhập & phiên người dùng và dữ liệu đúng tài khoản \\
UC02 & Tiếp nhận ảnh & chụp ảnh hoặc chọn từ thư viện & ảnh được kiểm tra và hiển thị xem trước \\
UC03 & Nhận diện trang phục & ảnh có một trang phục chính & nhãn dự đoán và độ tin cậy \\
UC04 & Quản lý tủ đồ & thêm, xem, sửa hoặc xoá món đồ & tủ đồ được cập nhật theo tài khoản \\
UC05 & Gợi ý phối đồ & món đồ đầu vào và dữ liệu tủ đồ & tối đa ba phương án không trùng \\
UC06 & Lưu và yêu thích & trang phục, phối đồ hoặc kết quả & nội dung được lưu để xem lại \\
UC07 & Thử đồ 2D & ảnh người dùng và ảnh trang phục & lớp ảnh có thể di chuyển, co giãn và xoay \\
\bottomrule
\end{tabularx}
\end{table}

Luồng xử lý chi tiết của bảy use case được minh hoạ lần lượt qua chức năng
xác thực ở Hình~\ref{fig:uc01-xac-thuc}, tiếp nhận ảnh ở
Hình~\ref{fig:uc02-tiep-nhan-anh}, nhận diện trang phục ở
Hình~\ref{fig:uc03-nhan-dien}, quản lý tủ đồ ở Hình~\ref{fig:uc04-tu-do},
gợi ý phối đồ ở Hình~\ref{fig:uc05-goi-y}, lưu nội dung và yêu thích ở
Hình~\ref{fig:uc06-luu-yeu-thich}, và thử trang phục 2D ở
Hình~\ref{fig:uc07-thu-do-2d}. Các sơ đồ cụ thể hoá những thao tác đã tổng
hợp trong Bảng~\ref{tab:usecase} mà không bổ sung chức năng ngoài phạm vi
đã xác định.

\begin{figure}[H]
\centering
\includegraphics[width=0.98\textwidth]{use-cases/uc01-xac-thuc.png}
\caption{Use case UC01: đăng ký và đăng nhập.}
\label{fig:uc01-xac-thuc}
\end{figure}

\begin{figure}[H]
\centering
\includegraphics[width=0.98\textwidth]{use-cases/uc02-tiep-nhan-anh.png}
\caption{Use case UC02: chọn hoặc chụp ảnh.}
\label{fig:uc02-tiep-nhan-anh}
\end{figure}

\begin{figure}[H]
\centering
\includegraphics[width=0.98\textwidth]{use-cases/uc03-nhan-dien.png}
\caption{Use case UC03: nhận diện trang phục.}
\label{fig:uc03-nhan-dien}
\end{figure}

\begin{figure}[H]
\centering
\includegraphics[width=0.98\textwidth]{use-cases/uc04-tu-do.png}
\caption{Use case UC04: quản lý tủ đồ.}
\label{fig:uc04-tu-do}
\end{figure}

\begin{figure}[H]
\centering
\includegraphics[width=0.98\textwidth]{use-cases/uc05-goi-y.png}
\caption{Use case UC05: gợi ý phối đồ.}
\label{fig:uc05-goi-y}
\end{figure}

\begin{figure}[H]
\centering
\includegraphics[width=0.98\textwidth]{use-cases/uc06-luu-yeu-thich.png}
\caption{Use case UC06: lưu nội dung và yêu thích.}
\label{fig:uc06-luu-yeu-thich}
\end{figure}

\begin{figure}[H]
\centering
\includegraphics[width=0.98\textwidth]{use-cases/uc07-thu-do-2d.png}
\caption{Use case UC07: thử trang phục 2D.}
\label{fig:uc07-thu-do-2d}
\end{figure}

Hệ thống hiện chỉ có một vai trò nghiệp vụ là người dùng đã xác thực. Firestore
cho phép đối chiếu \code{request.auth.uid} với định danh trên đường dẫn để giới
hạn người dùng vào dữ liệu của chính họ~\cite{FirestoreRules}. Quy tắc phiên bản
2 của dự án cho phép người dùng đọc, tạo, cập nhật và xoá đúng tài
liệu \code{users/\{userId\}} khi \code{request.auth.uid=userId}; quyền đọc/ghi
subcollection \code{wardrobe/\{itemId\}} cũng dùng cùng điều kiện. Người chưa
đăng nhập, truy cập UID khác hoặc truy cập đường dẫn không được khai báo đều
bị từ chối mặc định. Quy tắc chưa kiểm tra tên trường, kiểu dữ liệu hay giới
hạn nội dung, vì vậy người sở hữu vẫn có thể ghi mọi trường vào tài liệu của
mình; đây là điểm cần tăng cường trước khi triển khai thực tế.

Endpoint phân tích trang phục nhận multipart gồm ảnh JPG/JPEG/PNG tối đa
10 MB, hồ sơ, bối cảnh và ngưỡng confidence; phản hồi gồm các đối tượng
nhận diện và đánh giá Gemini nếu dịch vụ khả dụng. WebSocket thử đồ nhận cấu
hình camera trước, sau đó nhận JPEG nhị phân tối đa 3 MB và trả landmark/điểm
neo. Tiêu chí nghiệm thu là dữ liệu hợp lệ phải trả đúng schema; ảnh rỗng, sai
định dạng hoặc quá kích thước phải bị từ chối rõ ràng; tài khoản A không được
đọc/ghi dữ liệu của tài khoản B; ảnh và frame camera không được lưu tại backend.

\section{Thiết kế kiến trúc hệ thống}

Kiến trúc logic theo phạm vi đã mô tả được trình bày ở
Hình~\ref{fig:kientruc}. Sơ đồ chỉ thể hiện trách nhiệm của các khối; tên
dịch vụ, endpoint và vị trí triển khai phải được cập nhật theo mã nguồn.

\begin{figure}[H]
\centering
\resizebox{\textwidth}{!}{%
\begin{tikzpicture}[node distance=11mm and 14mm]
  \node[objbox] (user) {Người dùng\\Android};
  \node[clsbox, right=of user] (app) {Ứng dụng\\Flutter};
  \node[clsbox, right=of app] (api) {Dịch vụ\\API};
  \node[clsbox, right=of api] (ai) {YOLO, Gemini\\và MediaPipe};
  \node[objbox, right=of ai] (rec) {Kết quả phân tích,\\gợi ý và tư thế};
  \node[membox, below=of app] (overlay) {Ghép ảnh theo tư thế\\trên thiết bị};
  \node[membox, below=of api] (firebase) {Firebase Auth/Firestore\\tài khoản, hồ sơ, tủ đồ};
  \draw[flowarr] (user) -- (app);
  \draw[flowarr] (app) -- (api);
  \draw[flowarr] (api) -- (ai);
  \draw[flowarr] (ai) -- (rec);
  \draw[flowarr] (rec.south) -- ++(0,-8mm) -| (app.south);
  \draw[flowarr] (app) -- (overlay);
  \draw[flowarr] (app) -- (firebase);
\end{tikzpicture}}
\caption{Kiến trúc logic của hệ thống gợi ý phối đồ.}
\label{fig:kientruc}
\end{figure}

Ứng dụng Flutter chịu trách nhiệm tương tác, nhận ảnh, gọi API và hiển thị
kết quả. FastAPI là ranh giới với YOLO, Gemini, Open-Meteo và MediaPipe.
Flutter kết nối trực tiếp Firebase Authentication/Firestore; backend không
truy cập Firebase. Firestore đang lưu tài khoản và hồ sơ cơ thể, đồng thời
rules đã dành đường dẫn \code{users/\{uid\}/wardrobe/\{itemId\}} cho tủ đồ,
nhưng mã Flutter hiện chưa có luồng đọc/ghi subcollection này. Lớp trang phục
được Flutter ghép theo điểm neo tư thế và cho phép điều chỉnh thủ công.

\section{Thiết kế dữ liệu}

Từ mã nguồn ứng dụng, các nhóm dữ liệu thực tế được tổng hợp trong
Bảng~\ref{tab:dulieu}. Firebase dùng collection \code{users}; lịch sử hỏi
đáp và tuỳ chọn phân tích được lưu cục bộ theo định danh người dùng.

\begin{table}[H]
\centering
\caption{Các nhóm dữ liệu logic của hệ thống.}
\label{tab:dulieu}
\small
\begin{tabularx}{\textwidth}{>{\bfseries\color{daunavy}}l >{\raggedright\arraybackslash}X >{\raggedright\arraybackslash}X}
\toprule
\rowcolor{gray100}\bfseries\color{daunavy}Nhóm dữ liệu & \bfseries\color{daumaroon}Thông tin chính & \bfseries\color{daunavy}Mục đích \\
\midrule
Tài khoản & uid, họ tên, email, thời điểm tạo/cập nhật & xác thực và phân tách dữ liệu theo người dùng \\
Hồ sơ cơ thể & tên hồ sơ, chiều cao, cân nặng, giới tính tự khai, hồ sơ đang chọn & cá nhân hoá phân tích, gợi ý và hỏi đáp \\
Tuỳ chọn phân tích & chế độ thời tiết, mùa, nhiệt độ, phong cách, hoàn cảnh & khôi phục bối cảnh gần nhất theo tài khoản \\
Lịch sử hỏi đáp & vai trò, nội dung và thứ tự lượt chat & giữ hội thoại cục bộ và gửi tối đa sáu lượt gần nhất cho Gemini \\
Tủ đồ (đã có rules) & \code{users/\{uid\}/}\newline\code{wardrobe/\{itemId\}}; cấu trúc trường chưa xác minh & giới hạn quyền theo chủ tài khoản; mã Flutter hiện chưa sử dụng \\
Kết quả phân tích & lớp gốc, nhóm chuẩn hoá, màu, độ tin cậy, hộp và các điểm đánh giá & hiển thị kết quả YOLO/Gemini; không lưu ảnh ở backend \\
\bottomrule
\end{tabularx}
\end{table}

Ngoài dữ liệu nghiệp vụ, thành phần YOLO sử dụng hai nguồn ảnh. Bộ
\emph{Clothes Detection v1} xuất từ Roboflow gồm 11.926 ảnh, được chia thành
8.348 ảnh train, 2.385 ảnh validation và 1.193 ảnh test. File cấu hình công
bố giấy phép CC BY 4.0 và 21 lớp theo đúng thứ tự:
bag, belt, boat, coat, dress, glasses, hat, jacket, jumper, pants, pocket,
shirt, shoe, shoesClassic, shoesSports, shorts, skirt, sunglass, trousers,
tshirt và zipper.

Nguồn thứ hai là DeepFashion In-shop với 52.712 ảnh và hộp giới hạn
~\cite{DeepFashion2016}. Notebook V3 chuyển thành công 52.711 ảnh, bỏ qua một
ảnh lỗi, đồng thời ánh xạ 17 nhóm thư mục nguồn về taxonomy 21 lớp. Hai ánh
xạ gần đúng là \code{Rompers\_Jumpsuits} sang dress và \code{Suiting} sang
jacket. DeepFashion chỉ cho phép sử dụng phi thương mại trong nghiên cứu và
giáo dục theo README đi kèm dữ liệu.

Dữ liệu DeepFashion được chia theo \code{item\_id} với tỉ lệ 80/10/10 và
seed 42 để tránh rò rỉ nhiều góc chụp của cùng sản phẩm. Ảnh train của bộ dữ
liệu cũ được lặp hai lần trước khi trộn. Quy mô dữ liệu V3 sau khi chuẩn bị
được trình bày ở Bảng~\ref{tab:dulieu-yolo}.

\begin{table}[H]
\centering
\caption{Quy mô dữ liệu phát hiện trang phục ở phiên bản V3.}
\label{tab:dulieu-yolo}
\small
\begin{tabularx}{\textwidth}{>{\bfseries\color{daunavy}}l >{\raggedleft\arraybackslash}X >{\raggedleft\arraybackslash}X >{\raggedright\arraybackslash}p{6.0cm}}
\toprule
\rowcolor{gray100}\bfseries\color{daunavy}Tập & \bfseries\color{daumaroon}Ảnh & \bfseries\color{daumaroon}Đối tượng & \bfseries\color{daunavy}Ghi chú \\
\midrule
Train & 59.463 & 69.715 & gồm 16.696 lượt ảnh cũ sau khi lặp và 42.767 ảnh DeepFashion \\
Validation & 6.902 & 9.275 & gồm 2.385 ảnh cũ và 4.517 ảnh DeepFashion \\
Test & 6.620 & 7.821 & gồm 1.193 ảnh cũ và 5.427 ảnh DeepFashion \\
\midrule
Tổng & 72.985 & 86.811 & số ảnh train bao gồm lượt lặp có chủ đích \\
\bottomrule
\end{tabularx}
\end{table}

Nhật ký Ultralytics thực tế chấp nhận 59.315 ảnh train và 6.882 ảnh
validation; phần còn lại bị bỏ qua chủ yếu do nhãn cũ trộn hàng detection và
segmentation hoặc có toạ độ ngoài khoảng $[0,1]$. Đây là vấn đề chất lượng dữ
liệu cần được xử lý trước khi huấn luyện lại. Mã Flutter xác nhận hồ sơ tài
khoản và mảng \code{bodyProfiles} nằm trong \code{users/\{uid\}}; ảnh phân
tích, frame camera và kết quả YOLO không được ghi vào Firestore.

\section{Thiết kế mô hình hoặc cơ sở tri thức}

Thành phần nhận diện là bộ phát hiện đối tượng YOLO. Chuỗi xử lý gồm kiểm tra
ảnh, đưa ảnh về kích thước 640 pixel, suy luận hộp giới hạn và điểm lớp, lọc
kết quả theo ngưỡng của dịch vụ rồi trả nhãn, độ tin cậy và vị trí trang phục.
Các notebook và checkpoint cho thấy ba vòng lặp huấn luyện trong
Bảng~\ref{tab:cauhinh-yolo}; đây là các phiên bản fine-tune của cùng họ mô
hình chứ không phải ba kiến trúc độc lập.

\begin{table}[H]
\centering
\caption{Cấu hình các phiên bản huấn luyện YOLO.}
\label{tab:cauhinh-yolo}
\footnotesize
\begin{tabularx}{\textwidth}{>{\bfseries\color{daunavy}}l >{\raggedright\arraybackslash}p{2.4cm} >{\raggedright\arraybackslash}X >{\raggedright\arraybackslash}p{4.2cm}}
\toprule
\rowcolor{gray100}\bfseries\color{daunavy}Phiên bản & \bfseries\color{daumaroon}Khởi tạo & \bfseries\color{daunavy}Dữ liệu & \bfseries\color{daumaroon}Cấu hình chính \\
\midrule
V1 & \code{yolo26n.pt} & Roboflow 11.926 ảnh & tối đa 60 epoch, batch 16, ảnh 640, optimizer tự động \\
V2 & \code{best.pt} & dữ liệu cũ và 6.500 ảnh DeepFashion được chọn & tối đa 40 epoch, patience 10, AdamW, $lr_0=10^{-4}$, cosine LR \\
V3 & checkpoint V1/V3 & dữ liệu ở Bảng~\ref{tab:dulieu-yolo} & 30 epoch đã ghi nhận và 5 epoch tiếp tục, AdamW, batch 16, ảnh 640, seed 42 \\
\bottomrule
\end{tabularx}
\end{table}

Ở V3, lần chạy chính dùng $lr_0=10^{-4}$, weight decay 0,0005, cosine LR,
mosaic 0,5 và AMP. Sau khi phiên chạy dừng ở epoch 30, notebook nạp
\code{last.pt} và chạy tiếp 5 epoch với $lr_0=10^{-5}$, tắt AMP và lưu
checkpoint \code{best.pt}. Các tệp kết quả ghi nhận box loss, classification
loss và L1 loss cho cả train và validation.

Thành phần phân tích gửi metadata nhận diện, hồ sơ và bối cảnh đến Gemini để
nhận các điểm màu sắc, mùa/nhiệt độ, tỉ lệ cơ thể, phong cách và hoàn cảnh.
Điểm tổng dùng trọng số tương ứng 0,30; 0,30; 0,20; 0,10 và 0,10; tiêu chí
thiếu dữ liệu được bỏ khỏi mẫu số thay vì tính bằng 0. Luồng Gợi ý độc lập
với ảnh và YOLO, trả một nhận định, một lý do và từ một đến ba điều chỉnh.
Khi Gemini lỗi, hệ thống giữ kết quả YOLO và không tạo đánh giá giả.

Thành phần thử đồ nhận JPEG camera, chạy MediaPipe Pose ở chế độ
\code{VIDEO}, trả landmark và điểm neo qua WebSocket. Flutter ánh xạ toạ độ
camera, làm mượt tư thế rồi đặt lớp ảnh theo vai, hông hoặc bàn chân; người
dùng có thể tinh chỉnh vị trí, tỉ lệ và góc xoay. Phương pháp này không suy
diễn chính xác hình dạng ba chiều của cơ thể hoặc đặc tính vật lý của vải.

API mặc định nạp \code{backend/models/best.pt}; endpoint phân tích dùng
confidence 0,25 và IoU 0,45. Trên giao diện, người dùng có thể điều chỉnh
confidence trong khoảng 0,10--0,90. Các thông tin này được lấy trực tiếp từ
mã backend và mã Flutter.
